\roadmapSection{8}{MATHEMATICS}{30 Min/Day, Every Day}

{\small\itshape\color{arligray} 30 minutes every day. Use `Mastering the Fundamentals of Mathematics' as base book. No exceptions, even if tired.}

% ── 8.1 Common Mathematics ───────────────────────────────────────────────────
\roadmapSubSection{8.1}{Common Mathematics}{}

\checkitem{Algebra: rearrange equations, solve for unknown --- Ohm's Law, voltage dividers}{}
\checkitem{Scientific notation: 4.7$\mu$F = 4.7 $\times$ 10$^{-6}$ F --- essential for reading component values}{}
\checkitem{Unit conversions: mV/V/kV, mA/A, nF/$\mu$F/mF, kHz/MHz/GHz}{}
\checkitem{Ratios and percentages: efficiency = Pout/Pin $\times$ 100, duty cycle = ton/T $\times$ 100}{}
\checkitem{Basic statistics: mean, standard deviation --- for analyzing measurement data}{}

% ── 8.2 Mathematics for Electronics and Repair ───────────────────────────────
\roadmapSubSection{8.2}{Mathematics for Electronics and Repair}{}

\checkitem{Ohm's Law: V = I $\times$ R --- apply to every circuit you measure}{}
\checkitem{Power: P = V $\times$ I = V$^{2}$/R = I$^{2}$ $\times$ R --- understand what each form is useful for}{}
\checkitem{KVL and KCL: set up and solve simple loop equations for voltage dividers}{}
\checkitem{RC time constant: $\tau$ = R $\times$ C --- understand for capacitor charge/discharge}{}
\checkitem{Decibels (dB): gain in dB = 20 $\times$ log$_{10}$(Vout/Vin) --- used for signal attenuation}{}
\checkitem{Frequency: f = 1/T --- relate period to frequency for oscilloscope readings}{}

% ── 8.3 Mathematics for Digital Logic and CS ─────────────────────────────────
\roadmapSubSection{8.3}{Mathematics for Digital Logic and CS}{}

\checkitem{Binary: convert decimal $\leftrightarrow$ binary fluently without calculator}{}
\checkitem{Hexadecimal: convert binary $\leftrightarrow$ hex fluently --- 4 bits = 1 hex digit}{}
\checkitem{Two's complement: how negative integers are stored in binary}{}
\checkitem{Boolean algebra: AND, OR, NOT, NAND, NOR, XOR --- truth tables, De Morgan's theorem}{}
\checkitem{Bitwise operations: AND for masking, OR for setting, XOR for toggling, shifts for multiply/divide by 2}{}

% ── 8.4 Mathematics for Programming and Algorithms ───────────────────────────
\roadmapSubSection{8.4}{Mathematics for Programming and Algorithms}{}

\checkitem{Integer overflow: understand wrapping behavior in C --- important for security}{}
\checkitem{Modular arithmetic: x mod n --- used in ring buffers, hash functions}{}
\checkitem{Logarithms: understand O(log n) complexity --- binary search, balanced trees}{}
\checkitem{Big-O notation: O(1), O(n), O(n log n), O(n$^{2}$) --- analyze algorithm cost}{}
\checkitem{Amortized analysis: why dynamic array doubling is O(1) amortized}{}

% ── 8.5 Mathematics for PCB Design and Signals ───────────────────────────────
\roadmapSubSection{8.5}{Mathematics for PCB Design and Signals}{}

\checkitem{Trigonometry: sin and cos represent AC signal shape --- no derivation needed}{}
\checkitem{RC filter cutoff: $f_c$ = 1/(2$\pi$RC) --- plug in values, understand what changes}{}
\checkitem{LC resonance: $f$ = 1/(2$\pi\sqrt{LC}$) --- understand for RF filter design}{}
\checkitem{Impedance: Z = $\sqrt{R^2 + X^2}$ --- understand for 50$\Omega$ trace impedance control}{}
\checkitem{dB scale for antenna gain and signal path loss}{}

\newpage
% ── 8.6 Mathematics for Reverse Engineering and Security ─────────────────────
\roadmapSubSection{8.6}{Mathematics for Reverse Engineering and Security}{}

\checkitem{Modular exponentiation: $a^b \bmod n$ --- the core of RSA encryption}{}
\checkitem{XOR properties: A XOR A = 0, A XOR 0 = A --- why XOR is used in encryption}{}
\checkitem{Hash function properties: one-way, collision resistance --- used in firmware signing}{}
\checkitem{Finite fields GF($2^8$): used in AES --- concept only, not derivation}{}
\checkitem{CRC calculation: polynomial division in GF(2) --- how to verify firmware integrity}{}
